\subsection{The archimedean reflection extension: exact matrices, boundary
map, and domain-sensitive index comparison}

The next source locus, lines~573--624, defines one genuine local object:
the reflection crossed product $C_0(\mathbb R)\rtimes(\mathbb Z/2)$.  It does
not define the larger multi-place algebra mentioned above.  Lines~626--648
then write a scalar ``global component formula'', but do not provide a global
algebra, ideal, exact sequence, codiagonal map, or component-to-global
functor.  We therefore first compute the local crossed product exactly, then
separate three inequivalent domains for the displayed differential
expression, and only after that type the additional map required by the
global formula.

Let $\sigma(x)=-x$, let $u$ denote the unitary implementing $\sigma$, and put
\[
 S=\begin{pmatrix}0&1\\1&0\end{pmatrix},\qquad
 H=\frac1{\sqrt2}\begin{pmatrix}1&1\\1&-1\end{pmatrix},\qquad
 D_2=\left\{\begin{pmatrix}z_+&0\\0&z_-\end{pmatrix}:z_\pm\in\mathbb C\right\}.
\]
The matrix $H$ is retained: it is the explicit change from the orbit basis
$(r,-r)$ to the two character coordinates of $\mathbb Z/2$.

\begin{theorem}[Reflection crossed-product boundary]
\label{thm:primitive-archimedean-boundary}
There is a $*$-isomorphism
\[
 \Phi:C_0(\mathbb R)\rtimes_\sigma(\mathbb Z/2)
 \longrightarrow
 \mathcal A_\infty
 :=\{G\in C_0([0,\infty),M_2(\mathbb C)):G(0)\in D_2\}
\]
whose value on $f_0+f_1u$ is
\begin{equation}
 \Phi(f_0+f_1u)(r)
 =H\begin{pmatrix}
       f_0(r)&f_1(r)\\
       f_1(-r)&f_0(-r)
    \end{pmatrix}H,\qquad r\geq0.
 \label{eq:primitive-archimedean-Phi}
\end{equation}
Evaluation at the fixed point gives the exact sequence
\begin{equation}
 0\longrightarrow C_0((0,\infty),M_2)
 \longrightarrow\mathcal A_\infty
 \xrightarrow{\operatorname{ev}_0}D_2\longrightarrow0.
 \label{eq:primitive-archimedean-extension}
\end{equation}
Under $D_2\cong\mathbb C\oplus\mathbb C$, the projections
\[
 p_+=\frac{1+u}{2},\qquad p_-=\frac{1-u}{2}
\]
map to $e_{11}$ and $e_{22}$, respectively.

Identify
\[
 K_1(C_0((0,\infty),M_2))\cong\mathbb Z
\]
by declaring the loop
\[
 \beta(r)=\exp\!\left(2\pi i\frac{r}{1+r}\right),\qquad
 r\in[0,\infty],
\]
to have winding $+1$.  Then the exponential connecting map of
\eqref{eq:primitive-archimedean-extension} is
\begin{equation}
 \partial_\infty:\mathbb Z[p_+]\oplus\mathbb Z[p_-]\longrightarrow\mathbb Z,
 \qquad
 \partial_\infty(a,b)=-(a+b).
 \label{eq:primitive-archimedean-boundary-map}
\end{equation}
Consequently
\[
 \partial_\infty[p_+]=\partial_\infty[p_-]=-1,\qquad
 \partial_\infty([p_+]+[p_-])=-2,\qquad
 \ker\partial_\infty=\mathbb Z(1,-1).
\]
\end{theorem}

\begin{proof}
In the orbit basis of $\{r,-r\}$, the covariant matrix of the algebraic
crossed-product element $f_0+f_1u$ is
\[
 M_f(r)=
 \begin{pmatrix}
   f_0(r)&f_1(r)\\
   f_1(-r)&f_0(-r)
 \end{pmatrix}.
\]
At $r=0$ it has the form $\bigl(\begin{smallmatrix}a&b\\b&a\end{smallmatrix}\bigr)$,
and direct multiplication gives
\[
 H\begin{pmatrix}a&b\\b&a\end{pmatrix}H
 =\begin{pmatrix}a+b&0\\0&a-b\end{pmatrix}.
\]
Thus \eqref{eq:primitive-archimedean-Phi} lands in $\mathcal A_\infty$.

Conversely, for $G\in\mathcal A_\infty$ put $M(r)=HG(r)H$.  For $r>0$
define, without identifying any of the four coordinates,
\begin{equation}
 \begin{aligned}
 f_0(r)&=M_{11}(r),& f_0(-r)&=M_{22}(r),\\
 f_1(r)&=M_{12}(r),& f_1(-r)&=M_{21}(r).
 \end{aligned}
 \label{eq:primitive-archimedean-inverse-coordinates}
\end{equation}
Because $G(0)$ is diagonal, $M(0)$ has the form
$\bigl(\begin{smallmatrix}a&b\\b&a\end{smallmatrix}\bigr)$.  Hence the two
definitions of $f_0(0)$ agree, the two definitions of $f_1(0)$ agree, and all
four half-line coordinates glue continuously.  Since $G$ vanishes at
infinity, $f_0,f_1\in C_0(\mathbb R)$.  Formula
\eqref{eq:primitive-archimedean-inverse-coordinates} is therefore an explicit
inverse to $\Phi$.

For completeness at the norm level, Mesland--\c Seng\"un define the full and
reduced crossed-product norms and the regular Hilbert-module representation
in their Sections~1.3.16--1.3.20
\cite[pp.~12--13]{meslandSengun2025}.  The group $\mathbb Z/2$ is amenable by
the explicit invariant mean $m(h)=\frac12(h(0)+h(1))$.  Applied to the
semidirect-product Fell bundle of this action, Exel's approximation theorem
therefore identifies the full and reduced cross-sectional algebras
\cite[Theorems~20.6--20.7, pp.~148--149]{exel2017}.  The regular norm here is
precisely $\sup_{r\geq0}\|M_f(r)\|$.  Thus the mutually inverse coordinate maps
above extend isometrically to the stated $C^*$-algebras.

Evaluation at zero is onto $D_2$: for any $(z_+,z_-)$, choose
$h\in C_0([0,\infty))$ with $h(0)=1$ and use
$G(r)=h(r)\operatorname{diag}(z_+,z_-)$.  Its kernel consists exactly of the
matrix functions vanishing at both $0$ and infinity, namely
$C_0((0,\infty),M_2)$.  Moreover
\[
 HSH=\operatorname{diag}(1,-1),
\]
so $p_+$ and $p_-$ become $e_{11}$ and $e_{22}$.

Take the two self-adjoint lifts
\[
 x_+(r)=\begin{pmatrix}(1+r)^{-1}&0\\0&0\end{pmatrix},\qquad
 x_-(r)=\begin{pmatrix}0&0\\0&(1+r)^{-1}\end{pmatrix}.
\]
The exponential boundary formula gives
\[
 \partial_\infty[p_\pm]=[\exp(2\pi ix_\pm)]
\]
\cite[Theorem~9.3.1 and Section~9.3.2]{blackadar1986}.  In either case the
determinant of the nontrivial block is
\[
 \exp\!\left(\frac{2\pi i}{1+r}\right)
 =\exp\!\left(-2\pi i\frac r{1+r}\right)=\beta(r)^{-1}.
\]
It starts and ends at $1$ and winds once clockwise, hence represents $-1$ in
the declared generator.  Additivity gives
\eqref{eq:primitive-archimedean-boundary-map}, including its displayed
kernel.  No sign has been normalized away.
\end{proof}

The boundary computation is therefore correct.  The source's subsequent
index comparison requires a separate domain audit.

\begin{proposition}[The displayed differential expression does not determine
an APS index]
\label{prop:primitive-archimedean-domain-audit}
Let
\[
 F_t=(1-2t)I_2,\qquad 0\leq t\leq1,\qquad
 \mathcal D_F=\frac d{dt}+F_t.
\]
Then the following three exact realizations must not be conflated.
\begin{enumerate}[label=(\alph*)]
 \item Extend $F$ continuously by $+I_2$ for $t\leq0$ and by $-I_2$ for
 $t\geq1$, and let
 \[
 D_{\mathbb R}:W^{1,2}(\mathbb R,\mathbb C^2)
 \longrightarrow L^2(\mathbb R,\mathbb C^2),
 \qquad D_{\mathbb R}=\partial_t+F_t.
 \]
 Then $\operatorname{Ind}(D_{\mathbb R})=-2$.
 \item On the compact interval, with domain
 $H^1([0,1],\mathbb C^2)$ and no endpoint condition,
 $\operatorname{Ind}(\mathcal D_F)=+2$.
 \item On the compact interval, with domain
 $\{v\in H^1:v(0)=v(1)=0\}$,
 $\operatorname{Ind}(\mathcal D_F)=-2$.
\end{enumerate}
Hence the local source's notation
$\operatorname{Ind}_{\mathrm{APS}}(\mathcal D_F)=2$ is underdetermined until
its endpoint projectors, operator domain, adjoint domain, and orientation
convention are supplied.  For realization~(a), the exact relation is
\[
 \partial_\infty([p_+]+[p_-])
 =\operatorname{Ind}(D_{\mathbb R})=-2,
\]
whereas realization~(b) gives
$\partial_\infty([p_+]+[p_-])=-\operatorname{Ind}(\mathcal D_F)$.
\end{proposition}

\begin{proof}
Doll, Schulz-Baldes, and Waterstraat count negative-to-positive crossings
positively and prove, for
\[
 D_H=\partial_t-H_t:W^{1,2}(\mathbb R,\mathbb C^N)\to
 L^2(\mathbb R,\mathbb C^N),
\]
that
\begin{equation}
 \operatorname{Ind}(D_H)=-\operatorname{Sf}(H)
 \label{eq:primitive-doll-index-spectral-flow}
\end{equation}
when the asymptotic limits are invertible
\cite[Definition~1.1.3 and Theorem~3.5.1, pp.~3 and 78--82]
{dollSchulzBaldesWaterstraat2023}.  Their complete proof identifies the kernel
and cokernel through the stable and unstable spaces and obtains
\[
 \operatorname{Ind}(D_H)
 =\dim E^u(H_{-\infty})-\dim E^u(H_{+\infty}).
\]

Here $D_{\mathbb R}=D_H$ for $H=-F$.  The path $F$ goes from $+I_2$ to
$-I_2$, so the retained signature formula gives
\[
 \operatorname{Sf}(F)=\frac12((-2)-2)=-2,
 \qquad
 \operatorname{Sf}(-F)=\frac12(2-(-2))=2.
\]
Equation~\eqref{eq:primitive-doll-index-spectral-flow} therefore yields
$\operatorname{Ind}(D_{\mathbb R})=-2$.

On $[0,1]$, every homogeneous solution is
\[
 v(t)=e^{-t+t^2}c,\qquad c\in\mathbb C^2.
\]
With no endpoint condition this gives a two-dimensional kernel.  Variation of
constants with initial value $v(0)=0$ solves $\mathcal D_Fv=g$ for every
$g\in L^2$, so the operator is onto and its index is $2$.

With both endpoint values zero, the nonvanishing scalar
$e^{-t+t^2}$ forces $c=0$, hence the kernel is zero.  Integration by parts
shows that the adjoint is $-\partial_t+F_t$ with no endpoint condition.  Its
homogeneous solutions are
\[
 w(t)=e^{t-t^2}d,\qquad d\in\mathbb C^2,
\]
so the cokernel has dimension two and the index is $-2$.  These two compact
interval domains and the whole-line Fredholm domain are distinct data.  Since
the source supplies none of the boundary projections or domains needed to
select an APS realization, its bare differential expression cannot select
one of the resulting indices.
\end{proof}

\begin{proposition}[The scalar three-component formula needs an additional
codiagonal]
\label{prop:primitive-global-codiagonal}
Let the already proved finite-prime boundary be
$\partial_p(k)=-k$, and let $\partial_\infty$ be
\eqref{eq:primitive-archimedean-boundary-map}.  The direct sum of the two
defined extensions has connecting map
\[
 \partial_p\oplus\partial_\infty:
 \mathbb Z^3\longrightarrow\mathbb Z^2,\qquad
 (k,n_+,n_-)\longmapsto(-k,-n_+-n_-).
\]
The source's scalar formula is not this direct-sum boundary.  It is the
composite with the additional codiagonal
\[
 \Sigma:\mathbb Z^2\to\mathbb Z,\qquad \Sigma(a,b)=a+b,
\]
namely
\[
 (\Sigma\circ(\partial_p\oplus\partial_\infty))(k,n_+,n_-)
 =-(k+n_++n_-).
\]
Its kernel has the explicit $A_2$ basis
\[
 b_1=(1,-1,0),\qquad b_2=(0,1,-1),\qquad
 (\langle b_i,b_j\rangle)_{i,j}
 =\begin{pmatrix}2&-1\\-1&2\end{pmatrix}.
\]
By contrast, the direct-sum boundary has the rank-one kernel
$\mathbb Z(0,1,-1)$.  Thus the $A_2$ calculation is a correct calculation for
the displayed scalar composite.  At this stage of the source audit it is not
yet the boundary map of a defined multi-place $C^*$-extension: that admission
requires an algebraic construction inducing $\Sigma$, and the audited source
supplies no such map.  The independent stable Busby construction in
Theorem~\ref{thm:global-place-stable-extension} below closes this construction
obligation after stabilization; it does not retroactively attribute the map to
the local source or identify the source's undefined unstabilized global algebra.
\end{proposition}

\begin{proof}
The direct-sum formula follows component by component from the two proved
connecting maps.  Applying $\Sigma$ gives the displayed row vector
$-(1,1,1)$.  Both $b_1$ and $b_2$ lie in its kernel, they are independent, and
their dot products give the displayed Cartan matrix.  The direct-sum map has
two independent rows, so its kernel is exactly
$\mathbb Z(0,1,-1)$.  This proves both kernel statements and isolates the
additional information carried by $\Sigma$.
\end{proof}

\paragraph{Executable archimedean check.}
The bounded replay is
\begin{center}
\footnotesize
\path{certificates/globalization_nte_archimedean_boundary/verify_archimedean_boundary.py}.
\end{center}
It checks the Hadamard coordinate change, fixed-point boundary form, character
projections, explicit lifts and winding, both spectral-flow signs, the three
domain-sensitive index counts, the direct-sum and codiagonal matrices, and the
$A_2$ Gram matrix.  It does not certify the crossed-product completion, the
six-term theorem, Fredholmness, or an APS boundary theorem; the exact standard
proofs and source dependencies are printed above.
